\begin{problem}[The affine cross]\label{prob:vi18-p12}
Analyze $X=\Spec k[x,y]/(xy)$: find its two generic component points, their closures, the origin, and principal affine opens $D(x)$ and $D(y)$.
\end{problem}

\ifdefined\FullProblemDossiers
\paragraph{Why this problem matters.}
This problem consolidates the reducible example aspect of affine geometry and turns a structural statement into a reusable calculation.

\paragraph{Useful ideas before starting.}
Use the quotient ring and the relation $xy=0$.

\paragraph{Guided hints.}
\begin{enumerate}
\item Use the quotient ring and the relation $xy=0$.
\item Reduce the statement to a ring-theoretic property whenever possible.
\item Translate the result back into topology, sheaves, or local rings so that all affine-scheme layers are visible.
\end{enumerate}

\begin{solution}
Set
\[
A=k[x,y]/(xy).
\]
The minimal prime ideals are $(x)$ and $(y)$, because $xy=0$ and
\[
A/(x)\cong k[y],\qquad A/(y)\cong k[x]
\]
are domains. They are the generic points of the two irreducible components. Their closures are
\[
\overline{\{(x)\}}=V(x),\qquad
\overline{\{(y)\}}=V(y),
\]
the two coordinate axes.

The origin corresponds to the maximal ideal $(x,y)$. It lies on both components, since $(x)\subset(x,y)$ and $(y)\subset(x,y)$.

On $D(x)$, the element $x$ is invertible. The relation $xy=0$ then forces $y=0$, so
\[
A_x\cong k[x,x^{-1}].
\]
Hence
\[
D(x)\cong\Spec k[x,x^{-1}],
\]
the $x$-axis with the origin removed. Symmetrically,
\[
D(y)\cong\Spec k[y,y^{-1}],
\]
the punctured $y$-axis. Since $xy=0$, $D(x)\cap D(y)=D(xy)=\varnothing$. This exhibits explicitly how the reducible affine cross separates away from the intersection point.
\end{solution}


\paragraph{What happened mathematically.}
The calculation illustrates that global affine geometry is controlled by commutative algebra once the structure sheaf is kept in view.

\paragraph{Extensions.}
Vary the ring by products, quotients, localizations, arithmetic examples, or nilpotent thickenings and repeat the same dictionary.

\paragraph{Provenance.}
\textbf{New canonical synthesis; no numbered legacy Problem is claimed.}
\else
\paragraph{Provenance.} \textbf{New canonical synthesis; no numbered legacy Problem is claimed.}
\fi
